\documentclass[11pt]{amsart}
\usepackage{amsmath,amssymb,amsthm,mathtools}
\usepackage[margin=1in]{geometry}
\usepackage{array}
\usepackage{hyperref}

\newtheorem{theorem}{Theorem}[section]
\newtheorem{proposition}[theorem]{Proposition}
\newtheorem{lemma}[theorem]{Lemma}
\newtheorem{corollary}[theorem]{Corollary}
\theoremstyle{definition}
\newtheorem{definition}[theorem]{Definition}
\theoremstyle{remark}
\newtheorem{remark}[theorem]{Remark}

\newcommand{\RH}{\mathrm{RH}}
\renewcommand{\Re}{\operatorname{Re}}
\renewcommand{\Im}{\operatorname{Im}}
\newcommand{\half}{\tfrac12}
\newcommand{\fhat}{\widehat f}
\newcommand{\ghat}{\widehat g}
\newcommand{\HK}{\mathcal H_K}
\newcommand{\Fix}{\operatorname{Fix}}
\newcommand{\kV}{\kappa_{\mathcal V}}
\DeclareMathOperator{\Res}{Res}

\title[The localized Weil form on even test functions]{The localized Weil form on even test
functions:\\ off-line Lorentzian blocks and obstructions to positive-kernel realizations}
\author{David Alejandro Trejo Pizzo}
\thanks{Independent Researcher. \texttt{dtrejopizzo@gmail.com}}


\begin{document}
\begin{abstract}
For even real test functions $f\in C_c^\infty(\mathbb R)$ let
$Q(f,g)=\sum_\rho\varphi_f(\rho)\varphi_g(\rho)$, where $\varphi_f(s)=\int f(u)e^{(s-\frac12)u}\,du$
and $\rho$ runs over the nontrivial zeros of $\zeta$. Then $Q(f,f)$ is the zero side of Weil's
explicit formula at $f*f$, so by Weil's criterion the Riemann Hypothesis implies $Q\ge0$, and
it is equivalent to the corresponding positivity on all complex test functions. We give a
functional-analytic study of $Q$ on the even real space $\mathcal V$. Each on-line zero
contributes a square $\fhat(\gamma)^2\ge0$, and each off-line quartet $\{\half\pm b\pm i\gamma\}$
contributes a Lorentzian block $4\mu(u^2-v^2)$ with $u+iv=\fhat(\gamma-ib)$, $\mu$ the multiplicity. We prove that the
evaluation functionals at points with disjoint symmetry orbits are independent on $\mathcal V$;
hence, if there are $m<\infty$ off-line quartets, the off-line part of $Q$ has signature $(m,m)$
and the negative index $\kV$ of $Q$ satisfies $\kV\le m$; equality is left open. We prove three
elementary obstruction theorems: forms coming from a positive-definite reproducing kernel,
Cauchy-type Gram forms with parameters on one side of a line, and continuous non-degenerate
families of Hermitian forms on $\mathbb C^N$ cannot exhibit a negative direction that appears
exactly off the critical line. The resonance pairings of the modular surface are a sign-definite
instance. For the de Branges axioms we show that, for a realization of $Q$ in a Pontryagin space
of entire functions, the conjugation axiom holds under a natural compatibility condition, and
that a realization satisfying all three axioms is a space $H(E)$ whose negative index is at least
$\kV$ (Kaltenb\"ack--Woracek). We also note that $E_\xi(z)=\xi(\half+iz)$ satisfies
$E_\xi^\#=E_\xi$, so $H(E_\xi)=\{0\}$ and $E_\xi$ is not a structure function. Nothing here
proves or assumes the Riemann Hypothesis.
\end{abstract}
\maketitle

\tableofcontents
\newpage

\section{Introduction}

Let $\xi(s)=\half s(s-1)\pi^{-s/2}\Gamma(s/2)\zeta(s)$ be the completed Riemann zeta
function, entire of order one, with $\xi(s)=\xi(1-s)$, real on the real axis and on the critical
line, and with nontrivial zeros $\rho=\beta+i\gamma$ confined to the open strip $0<\beta<1$. The
Riemann Hypothesis ($\RH$) asserts $\beta=\half$ for all of them.

Weil's explicit formula~\cite{Weil} relates a sum over zeros to an archimedean term and a sum
over prime powers. For a suitable even real test function $h$ with transform $\varphi_h$
(Definition~\ref{def:setup}),
\begin{equation}\label{eq:weil-id}
\sum_\rho\varphi_h(\rho)\;=\;A_\infty(h)-P(h),
\end{equation}
with $A_\infty$ the archimedean (gamma-factor and pole) term and $P$ the sum over prime powers.
Weil's criterion~\cite{Weil,Bombieri} states that $\RH$ holds if and only if the left side of
\eqref{eq:weil-id} is nonnegative at $h=g*\tilde g$ for every $g\in C_c^\infty(\mathbb R)$, where
$\tilde g(u)=\overline{g(-u)}$. For even real $f$ one has $\tilde f=f$, and the zero side at $f*f$
is the symmetric bilinear form $Q(f,f)=\sum_\rho\varphi_f(\rho)^2$ (Remark~\ref{rem:choices}).
Thus $\RH$ implies $Q\ge0$ on the space $\mathcal V$ of even real test functions. We write
$\kV\in\{0,1,2,\dots,\infty\}$ for the negative index of $Q$ on $\mathcal V$, the supremum of the
dimensions of subspaces on which $Q$ is negative definite.

A long line of spectral approaches (Hilbert--P\'olya and its de Branges canonical-system form,
Connes' trace formula~\cite{Connes99}, Burnol's study of de Branges spaces attached to the
Fourier transform~\cite{Burnol}) seeks a Hilbert or Pontryagin space attached to the zeros in
which the positivity of such a form would be visible. Kre\u\i n-space and Pontryagin-index
formulations of the zero side of the Weil form were considered in earlier papers of this
series: a Kre\u\i n-space formulation on Paley--Wiener spaces in~\cite{P8}, and the negative
index of the full Weil form $A_\infty-P$ on functions of bounded support in~\cite{P11}; a
finite-dimensional study of the effect of an off-line pair on a Gram matrix of the Weil form is
in~\cite{P7}. The present paper studies the zero-side form $Q$ itself on $\mathcal V$.

The results are as follows.
\begin{itemize}
\item \textbf{Localization} (\S\ref{sec:setup}--\S\ref{sec:localization}). $Q$ is defined
directly by an absolutely convergent sum (Lemma~\ref{lem:conv}). On $\mathcal V$ it splits as a
sum of nonnegative squares $\fhat(\gamma)^2$ over on-line zeros and Lorentzian blocks
$4\mu(u^2-v^2)$, $u+iv=\fhat(\gamma-ib)$, over off-line quartets (Theorem~\ref{thm:loc}).
\item \textbf{Obstructions} (\S\ref{sec:modular}--\S\ref{sec:obstruction}). Forms realized by a
positive-definite kernel are positive semidefinite; Cauchy Gram forms $d_i\bar d_j/(z_i+\bar
z_j-c)$ with all $z_i$ on one side of $\Re z=c/2$ are definite; and a continuous family of
non-degenerate Hermitian forms on $\mathbb C^N$ has constant index. The resonance pairings of the
modular surface are definite for every configuration of zeros in the strip
(Theorem~\ref{thm:modular}).
\item \textbf{Mechanism and count} (\S\ref{sec:mechanism}--\S\ref{sec:regimes}). Evaluation
functionals at points with pairwise disjoint orbits $\{\pm z,\pm\bar z\}$ are independent on
$\mathcal V$ (Lemma~\ref{lem:indep}). If there are $m<\infty$ off-line quartets, the off-line part
of $Q$ has signature $(m,m)$ and $\kV\le m$ (Corollary~\ref{cor:count}). Whether $\kV=m$ is left
open; whether $m$ is finite is open (Proposition~\ref{prop:finiteness}).
\item \textbf{de Branges axioms} (\S\ref{sec:rigidity}). For a realization of $Q$ in a Pontryagin
space of entire functions, the conjugation axiom follows from a compatibility condition; if all
three axioms hold, the space is $H(E)$ for a structure function $E$ with $k(E)\ge\kV$ zeros in the
upper half-plane, by the theorem of Kaltenb\"ack and Woracek~\cite{KW} (Theorem~\ref{thm:rigid}).
\end{itemize}
Section~\ref{sec:dichotomy} records the symmetries of the zero set: the Klein four-group generated
by $s\mapsto\bar s$ and $s\mapsto1-s$ preserves the zeros unconditionally, while the statement
that $s\mapsto1-\bar s$ fixes every zero is $\RH$. Heuristic remarks on the size of negative
directions are labelled as such.

None of the results assumes $\RH$, and none of them proves it.

\section{The form \texorpdfstring{$Q$}{Q} and its convergence}\label{sec:setup}

We define $Q$ directly on a concrete test space.

\begin{definition}[Test space and transform]\label{def:setup}
Let $\mathcal V$ be the space of even real $f\in C_c^\infty(\mathbb R)$. For $f\in\mathcal V$
define the entire transform
\[
\varphi_f(s)\;=\;\int_{\mathbb R} f(u)\,e^{(s-\frac12)u}\,du,\qquad s\in\mathbb C .
\]
Since $f$ has compact support, $\varphi_f$ is entire of exponential type. Writing
$s=\half+iz$ we have $\varphi_f(\half+iz)=\int f(u)e^{izu}\,du=\fhat(z)$, the Fourier
transform of $f$ extended to $z\in\mathbb C$; because $f$ is even and real, $\fhat$ is even
and real on $\mathbb R$.
\end{definition}

\begin{definition}[The localized Weil form]\label{def:weil}
For $f,g\in\mathcal V$ set
\[
Q(f,g)\;=\;\sum_{\rho}\varphi_f(\rho)\,\varphi_g(\rho),
\]
the sum running over the nontrivial zeros $\rho$ of $\zeta$ with multiplicity.
\end{definition}

\begin{lemma}[Absolute convergence and symmetries]\label{lem:conv}
For $f,g\in\mathcal V$ the series in Definition~\ref{def:weil} converges absolutely, and
\[
\varphi_f(1-s)=\varphi_f(s),\qquad \varphi_f(\bar s)=\overline{\varphi_f(s)} .
\]
\end{lemma}
\begin{proof}
Evenness of $f$ gives $\varphi_f(1-s)=\int f(u)e^{(\frac12-s)u}du=\int f(-u)e^{(\frac12-s)u}
du=\int f(u)e^{(s-\frac12)u}du=\varphi_f(s)$; reality of $f$ gives
$\overline{\varphi_f(s)}=\int f(u)e^{(\bar s-\frac12)u}du=\varphi_f(\bar s)$. For
convergence, write $\rho=\beta+i\gamma$. With $f$ supported in $[-T_0,T_0]$,
$|\varphi_f(\rho)|\le\|f\|_1\,e^{T_0|\beta-\frac12|}\le\|f\|_1\,e^{T_0/2}$, while integrating by
parts twice against $e^{(\rho-\frac12)u}$ yields $|\varphi_f(\rho)|\le
\|f''\|_1e^{T_0/2}\,|\rho-\half|^{-2}$, uniformly for $0\le\beta\le1$. The Riemann--von
Mangoldt formula $\#\{\rho:0<\gamma\le T\}\sim\frac{T}{2\pi}\log\frac{T}{2\pi}$ makes
$\sum_\rho(1+|\gamma|)^{-2}<\infty$, hence $\sum_\rho|\varphi_f(\rho)\varphi_g(\rho)|<\infty$.
\end{proof}

The two symmetries say that the summand $\varphi_f(\rho)^2$ is invariant under
$\rho\mapsto1-\rho$ and is conjugated under $\rho\mapsto\bar\rho$; this is what couples the
four members of an off-line orbit.

\section{Localization of the negative directions in off-axis evaluation}\label{sec:localization}

The nontrivial zeros split into the on-line zeros $\rho=\half+i\gamma$ and the off-line
zeros, which by the functional equation and the reality of $\xi$ occur in \emph{Klein quartets}
\[
\{\rho,\ \bar\rho,\ 1-\rho,\ 1-\bar\rho\}=\{\tfrac12\pm b\pm i\gamma\},\qquad
b=\beta-\tfrac12\neq0 ,
\]
whose four members have a common multiplicity $\mu$. Since $\zeta$ has no zeros on $(0,1)$,
$\gamma\neq0$ for every nontrivial zero. For a point $z\in\mathbb C$ write
$E_z(f)=\fhat(z)$ for $f\in\mathcal V$.

\begin{theorem}[Localization of the negative directions]\label{thm:loc}
On $\mathcal V$ the Weil form decomposes as $Q=Q_{\mathrm{on}}+Q_{\mathrm{off}}$, both series
converging absolutely, where:
\begin{enumerate}
\item each on-line zero $\half+i\gamma$ contributes the nonnegative term
$\varphi_f(\half+i\gamma)^2=\fhat(\gamma)^2\ge0$, with $\fhat(\gamma)\in\mathbb R$;
\item each off-line quartet $\{\half\pm b\pm i\gamma\}$ of multiplicity $\mu$ contributes the
Lorentzian block
\[
4\mu\,\Re\!\big[\fhat(\gamma-ib)^2\big]\;=\;4\mu\,(u^2-v^2),\qquad
u=\Re\fhat(\gamma-ib),\ \ v=\Im\fhat(\gamma-ib),
\]
which depends on $f$ only through the off-axis evaluation $E_{\gamma-ib}(f)$.
\end{enumerate}
In particular $Q_{\mathrm{on}}\ge0$, and $\RH$ implies $Q_{\mathrm{off}}=0$ and $Q\ge0$ on
$\mathcal V$. Conversely, if the off-line zeros form $m$ quartets with $1\le m<\infty$, then
$Q_{\mathrm{off}}$ has signature $(m,m)$ on $\mathcal V$ \textnormal{(}Corollary~\ref{cor:count},
which uses Lemma~\ref{lem:indep} of Section~\ref{sec:mechanism}\textnormal{)}, so
$Q_{\mathrm{off}}\neq0$.
\end{theorem}
\begin{proof}
For an on-line zero, $\varphi_f(\half+i\gamma)=\fhat(\gamma)$ is real (Lemma~\ref{lem:conv}),
so its square is $\ge0$. For an off-line quartet, put $w=\varphi_f(\half+b+i\gamma)$; since
$\half+b+i\gamma=\half+i(\gamma-ib)$, $w=\fhat(\gamma-ib)$. By $\varphi_f(1-s)=\varphi_f(s)$ the
zero $\half-b-i\gamma$ gives $w$, and by $\varphi_f(\bar s)=\overline{\varphi_f(s)}$ the zeros
$\half+b-i\gamma$ and $\half-b+i\gamma$ give $\bar w$; the quartet contributes
$\mu(2w^2+2\bar w^2)=4\mu\Re[w^2]=4\mu(u^2-v^2)$. Both subseries converge absolutely by
Lemma~\ref{lem:conv}. The last assertion is Corollary~\ref{cor:count}(a).
\end{proof}

\begin{remark}[One block per quartet on the real space]\label{rem:two-blocks}
On the even real space $\mathcal V$ the identities $E_{-z}=E_z$ and
$E_{\bar z}=\overline{E_z}$ make the evaluations at the four points $\pm\gamma\pm ib$ of a quartet
expressible through the single complex functional $E_{\gamma-ib}$; a quartet therefore produces
exactly one Lorentzian block, with one negative and one positive direction. This is the source of
the bound $\kV\le m$ in Corollary~\ref{cor:count}.
\end{remark}

\begin{remark}[Relation with the Weil functional]\label{rem:choices}
For $f,g\in\mathcal V$ the convolution $f*g$ lies in $\mathcal V$ and
$\varphi_{f*g}=\varphi_f\varphi_g$. Hence $Q(f,g)=\sum_\rho\varphi_{f*g}(\rho)$ is the zero side
of \eqref{eq:weil-id} at $h=f*g$, and $Q(f,f)=A_\infty(f*f)-P(f*f)$. Since $\tilde f=f$ for
$f\in\mathcal V$, the values $Q(f,f)$ are the values of Weil's functional at $g*\tilde g$ with
$g=f$ even and real; Weil's criterion~\cite{Weil,Bombieri} uses all complex $g\in C_c^\infty$, and
we do not use the converse direction on $\mathcal V$ alone.
\end{remark}

Theorem~\ref{thm:loc} shows that the negative directions of $Q$, when they exist, come from the
off-axis point evaluations $E_{\gamma-ib}$ ($b\ne0$). On $L^2(\mathbb R)$ these evaluations are
unbounded, so a space carrying $Q$ as an inner product must support bounded off-axis evaluations.
The next two sections show which natural constructions are necessarily sign-definite.

\section{A definite test case: the modular-surface resonances}\label{sec:modular}

On the modular surface $X=SL_2(\mathbb Z)\backslash\mathbb H$ the zeros appear as scattering
resonances: the constant term of the Eisenstein series $E(z,s)$ is $y^s+\phi(s)y^{1-s}$ with
scattering coefficient
\[
\phi(s)=\frac{\Lambda(2s-1)}{\Lambda(2s)},\qquad
\Lambda(s)=\pi^{-s/2}\Gamma(s/2)\zeta(s)=\frac{2\,\xi(s)}{s(s-1)},
\]
see~\cite{Iwaniec}, whose poles in $0<\Re s<\half$ are the points
$s_\rho=\rho/2=\beta_\rho/2+i\gamma_\rho/2$ (since $\Lambda(2s)=0\iff2s=\rho$ there).

\begin{proposition}[The resonance weights]\label{prop:weights}
For a simple zero $\rho$, the residue of the scattering coefficient is
\[
R_\rho=\Res_{s=\rho/2}\phi(s)=\frac{\Lambda(\rho-1)}{2\,\Lambda'(\rho)}
=\frac{\rho}{\rho-2}\,\frac{\xi(\rho-1)}{2\,\xi'(\rho)},
\]
and, as $\gamma_\rho\to\infty$ with $\beta_\rho$ in a compact subset of $(0,1)$,
\[
|R_\rho|\sim\frac12\Big(\frac{\gamma_\rho}{2\pi}\Big)^{1-\beta_\rho}
\frac{|\zeta(2-\rho)|}{|\zeta'(\rho)|}.
\]
Numerically, $|R_\rho|$ lies between $0.51$ and $0.69$ for the first fifteen zeros (minimum
$0.515$ at the first zero, maximum $0.689$ at the $13$th), not monotonically, and equals $0.70$, $0.84$ and $1.04$ at the $50$th, $100$th and $300$th zeros.
\end{proposition}
\begin{proof}
The residue formula is immediate since $\Lambda(2s)$ has a simple zero at $s=\rho/2$ with
derivative $2\Lambda'(\rho)$; the second expression follows from $\xi(s)=\half s(s-1)\Lambda(s)$,
which gives $\xi(\rho-1)=\half(\rho-1)(\rho-2)\Lambda(\rho-1)$ and
$\xi'(\rho)=\half\rho(\rho-1)\Lambda'(\rho)$. For the asymptotic, since
$\Lambda'(\rho)=\pi^{-\rho/2}\Gamma(\rho/2)\zeta'(\rho)$ at a zero, write
\[
\frac{\Lambda(\rho-1)}{\Lambda'(\rho)}=\pi^{1/2}\,
\frac{\Gamma(\frac{\rho-1}{2})}{\Gamma(\frac{\rho}{2})}\,\frac{\zeta(\rho-1)}{\zeta'(\rho)} .
\]
By Stirling's formula $|\Gamma(\frac{\rho-1}{2})/\Gamma(\frac\rho2)|\sim(\gamma/2)^{-1/2}$, and by
the functional equation $|\zeta(\rho-1)|=|\chi(\rho-1)|\,|\zeta(2-\rho)|$ with
$|\chi(\sigma+it)|\sim(|t|/2\pi)^{1/2-\sigma}$, so
$|\zeta(\rho-1)|\sim(\gamma/2\pi)^{3/2-\beta}|\zeta(2-\rho)|$. Multiplying,
$|\Lambda(\rho-1)/\Lambda'(\rho)|\sim(\gamma/2\pi)^{-1/2}(\gamma/2\pi)^{3/2-\beta}|\zeta(2-\rho)|/|\zeta'(\rho)|$,
which gives the claim. The numerical values are computed by
\texttt{scripts/resonance\_weights.py} (mpmath, $\xi'$ by numerical differentiation); at the
$10$th, $100$th and $300$th zeros the asymptotic expression agrees with $|R_\rho|$ to three
digits.
\end{proof}

On the line $|\zeta(2-\rho)|=|\zeta(\frac32-i\gamma)|$ is bounded above and below, so
$|R_\rho|$ is of size $\gamma^{1/2}/|\zeta'(\rho)|$; no unconditional upper or lower bound for
$|R_\rho|$ is claimed.

\begin{theorem}[The resonance pairings are half-plane definite]\label{thm:modular}
Write $\rho=\beta_\rho+i\gamma_\rho$ with $0<\beta_\rho<1$, $s_\rho=\rho/2$, and let $R_\rho$ be
as in Proposition~\ref{prop:weights} for simple zeros $\rho$ (then $R_\rho\neq0$, since
$\Lambda(\rho-1)\neq0$ because $\Re(\rho-1)<0$).
\begin{enumerate}
\item \emph{(Sesquilinear.)} For any finite set of distinct simple zeros, the Hermitian matrix $G=DKD^{*}$,
$D=\mathrm{diag}(R_\rho)$, $K_{\rho\sigma}=(s_\rho+\overline{s_\sigma}-1)^{-1}$, is
negative definite, for every configuration of zeros in $0<\beta<1$.
\item \emph{(Bilinear.)} The bilinear pairing of $\rho$ with $\bar\rho$ defined by the same
Cauchy-type kernel,
\[
B(\rho,\bar\rho)=\frac{-R_\rho R_{\bar\rho}}{s_\rho+s_{\bar\rho}-1}
=\frac{|R_\rho|^2}{1-\beta_\rho},
\]
is positive for every $\beta_\rho\in(0,1)$, on-line and off-line alike.
\end{enumerate}
Consequently both pairings are sign-definite throughout the strip; neither has a negative
direction that appears exactly when $\beta_\rho\neq\half$.
\end{theorem}
\begin{proof}
(1) Set $w_\rho=\half-s_\rho$, so $\Re w_\rho=\tfrac{1-\beta_\rho}{2}>0$ because
$\beta_\rho<1$. Then
\[
K_{\rho\sigma}=\frac{1}{s_\rho+\overline{s_\sigma}-1}=\frac{-1}{w_\rho+\overline{w_\sigma}}
=-\int_0^\infty e^{-w_\rho t}\,\overline{e^{-w_\sigma t}}\,dt,
\]
the negative of the Szeg\H o kernel of the Hardy space of the right half-plane. Since the
$w_\rho$ are distinct, the functions $e^{-w_\rho t}$ are linearly independent in
$L^2(0,\infty)$, so $K$ is negative definite; conjugating by the invertible $D$ preserves this.
(2) Use $R_{\bar\rho}=\overline{R_\rho}$ (as $\phi(\bar s)=\overline{\phi(s)}$) and
$s_\rho+s_{\bar\rho}-1=\beta_\rho-1$.
\end{proof}

\begin{remark}
Matrices of the form $DKD^*$ arise from the Maass--Selberg relations for truncated Eisenstein
series~\cite[Ch.~6]{Iwaniec} and from Zagier's renormalization of Rankin--Selberg integrals of
functions not of rapid decay~\cite{Zagier}: taking residues at $s_1=s_\rho$, $s_2=s_\sigma$ in
the Maass--Selberg relation, the term involving $\phi(s_1)\overline{\phi(s_2)}$ contributes
$-R_\rho\overline{R_\sigma}\,T^{1-s_\rho-\overline{s_\sigma}}/(s_\rho+\overline{s_\sigma}-1)$.
Theorem~\ref{thm:modular} uses only the matrices themselves, and its conclusion holds in either
sign convention. The obstruction is that the kernels are of Cauchy type
$1/(s_\rho+\overline{s_\sigma}-c)$ with denominators of one sign throughout $0<\beta<1$, while a
negative direction caused by off-line zeros must appear exactly when $\beta\neq\half$.
Section~\ref{sec:obstruction} formulates this in general.
\end{remark}

\section{The obstruction theorems}\label{sec:obstruction}

We isolate three properties and show that each forces definiteness or a constant index. The
results are elementary.

\subsection{Positive kernels are semidefinite}
Let $\Omega$ be a set and $K:\Omega\times\Omega\to\mathbb C$ a positive-definite kernel; let
$\HK$ be its reproducing-kernel Hilbert space (Moore--Aronszajn~\cite{Aronszajn}), with
$K_z=K(\cdot,z)$ and $\langle f,K_z\rangle=f(z)$.

\begin{theorem}[Positive kernel $\Rightarrow$ semidefinite]\label{thm:A}
Any Hermitian form realized as $B(f,g)=\langle f,g\rangle_{\HK}$ for a positive-definite kernel
$K$ is positive semidefinite. In particular, if $J:\mathcal V\to\HK$ is real-linear and
$Q(f,g)=\Re\langle Jf,Jg\rangle_{\HK}$ for all $f,g\in\mathcal V$, then $\kV=0$.
\end{theorem}
\begin{proof}
For $f=\sum_i c_iK_{z_i}$, $\|f\|_{\HK}^2=\sum_{i,j}c_i\bar c_j K(z_i,z_j)\ge0$ by
positive-definiteness; such $f$ are dense~\cite{Aronszajn}, so $B\succeq0$. Then
$Q(f,f)=\|Jf\|^2\ge0$.
\end{proof}

Thus indefiniteness cannot be produced from a positive kernel: Bergman, Hardy, Fock, heat-kernel
and Bochner-positive convolution pairings, and de Branges spaces with a Hermite--Biehler structure
function, all give $\kV=0$ when used in this way.

\subsection{Fixed-sign Cauchy kernels are definite}
Theorem~\ref{thm:modular}(1) is the special case $c=1$ of the following.

\begin{theorem}[Fixed-sign Cauchy kernel $\Rightarrow$ definite]\label{thm:B}
Let $c\in\mathbb R$ and $S\subset\mathbb C$ with all $\Re z<c/2$ (resp.\ all $\Re z>c/2$).
Then $K(z,w)=(z+\bar w-c)^{-1}$ is negative definite (resp.\ positive definite) on finite subsets
of distinct points of $S$, and any Gram form $G=[d_i\bar d_j K(z_i,z_j)]$ with weights $d_i\ne0$
is negative (resp.\ positive) definite.
\end{theorem}
\begin{proof}
With $w_z=\tfrac c2-z$ (so $\Re w_z>0$ on $S$ in the first case),
$K(z,w)=-(w_z+\overline{w_w})^{-1}=-\int_0^\infty e^{-w_zt}\overline{e^{-w_wt}}\,dt$, which is
negative definite on distinct points as in the proof of Theorem~\ref{thm:modular}; the other case
is identical with $w_z=z-\tfrac c2$, and conjugation by $D=\mathrm{diag}(d_i)$ preserves
definiteness.
\end{proof}

The hypothesis is that the parameters lie in a half-plane not straddling $\Re z=c/2$. For the
modular resonances $s_\rho=\rho/2$ this holds because $0<\beta_\rho<1$ confines them to
$\Re s_\rho<\half$, whether or not $\RH$ holds.

\subsection{No continuous non-degenerate family detects the line}
The condition $\beta=\half$ has real codimension one; signature is locally constant on the
non-degenerate locus.

\begin{theorem}[Continuity $\Rightarrow$ no detection]\label{thm:C}
Let $\beta\mapsto Q_\beta$ be a continuous family of Hermitian forms on $\mathbb C^N$, defined for
$\beta$ near $\half$. If $Q_\beta$ is non-degenerate throughout a neighborhood of $\half$, its
negative index $\kappa(\beta)$ is constant there. In particular $Q_\beta$ cannot have
$\kappa(\half)=0$ and $\kappa(\beta)\ge1$ for $\beta\ne\half$ unless it degenerates at $\half$.
\end{theorem}
\begin{proof}
The eigenvalues of a Hermitian matrix depend continuously on its entries; $\kappa(\beta)$, the
number of negative ones, can change only where an eigenvalue crosses $0$, i.e.\ at a degeneracy.
On a connected non-degenerate neighborhood none crosses, so $\kappa$ is constant.
\end{proof}

For example, the bilinear pairing of Theorem~\ref{thm:modular}(2) equals
$|R_\rho|^2/(1-\beta)$, which is continuous and nonvanishing at $\beta=\half$, hence of constant
sign there. By contrast, the off-line displacement $b=\beta-\half$ that enters the evaluation
point $\gamma-ib$ of Theorem~\ref{thm:loc} vanishes at $\beta=\half$, where the block
$4\mu(u^2-v^2)$ degenerates to $4\mu u^2$.

\begin{corollary}[The excluded class and the two escape routes]\label{cor:class}
Suppose $\kV\ge1$ (which happens if $\RH$ fails with finitely many off-line quartets and
$\kV=m$, see Corollary~\ref{cor:count}). Then $Q$ on $\mathcal V$ is not of any of the following
forms: \emph{(i)} $\Re\langle Jf,Jg\rangle$ for a real-linear map $J$ into the
reproducing-kernel Hilbert space of a positive-definite kernel (Theorem~\ref{thm:A});
\emph{(ii)} a Gram form $d_i\bar d_j/(z_i+\bar z_j-c)$ with all $z_i$ on one side of
$\Re z=c/2$, in particular the Gram matrix of Theorem~\ref{thm:modular}(1)
(Theorem~\ref{thm:B}). Moreover \emph{(iii)} a continuous family of Hermitian forms on
$\mathbb C^N$ whose index is $0$ at $\beta=\half$ and positive for $\beta\neq\half$ must
degenerate at $\beta=\half$ (Theorem~\ref{thm:C}). Thus none of these three classes can carry the
negative directions of $Q$; the constructions considered below use instead an \emph{indefinite
(Pontryagin) kernel}, or a \emph{degeneration at $\beta=\half$}.
\end{corollary}

The second route is realized by the off-axis evaluation of Theorem~\ref{thm:loc}, which
degenerates as $b\to0$. The first leads to indefinite reproducing-kernel spaces, whose local
mechanism we now describe.

\section{The local mechanism and the count}\label{sec:mechanism}

\subsection{Bilinear, not sesquilinear}
For even real $f$, $\fhat$ is even and real on $\mathbb R$, so the entire extension of
$|\fhat|^2$ from $\mathbb R$ is the holomorphic square $\fhat(z)^2$, not $\fhat(z)\overline{\fhat(z)}$.
Hence
\begin{equation}\label{eq:bilinear}
Q(f,g)=\sum_\rho\fhat(z_\rho)\,\ghat(z_\rho),\qquad z_\rho=\tfrac{\rho-\frac12}{i},
\end{equation}
is a \emph{symmetric bilinear} form, not the sesquilinear $\sum\fhat(z_\rho)\overline{\ghat
(z_\rho)}$. The bilinear square of a complex number is indefinite.

\begin{proposition}[Sesquilinear vs.\ bilinear]\label{prop:vs}
The form $f\mapsto|E_z(f)|^2$ is $\ge0$ for every $z$. The form $f\mapsto\Re[E_z(f)^2]$ is
$\ge0$ where $E_z(f)$ is real for all $f\in\mathcal V$ (namely $z\in\mathbb R\cup i\mathbb R$),
and with $E_z(f)=u+iv$ it equals $u^2-v^2$ in general.
\end{proposition}

\subsection{The dichotomy theorems}
\begin{theorem}[Sesquilinear evaluation forms are positive]\label{thm:E}
In any reproducing-kernel Hilbert space with continuous evaluations $E_z(f)=\langle f,K_z
\rangle$, and any points $z_i$ and weights $c_i\ge0$, the form $f\mapsto\sum_i c_i|E_{z_i}(f)|^2$
is positive semidefinite.
\end{theorem}
\begin{proof}
$|E_{z_i}(f)|^2=|\langle f,K_{z_i}\rangle|^2\ge0$, and a nonnegative combination of
nonnegative forms is nonnegative.
\end{proof}

\begin{theorem}[Bilinear evaluation squares: signature $(1,1)$ off-axis]\label{thm:F}
On $\mathcal V$, the real quadratic form $f\mapsto\Re[E_z(f)^2]$ has signature $(1,0)$ if
$z\in\mathbb R\cup i\mathbb R$, and signature $(1,1)$ otherwise. Thus each off-axis evaluation
contributes exactly one negative direction.
\end{theorem}
\begin{proof}
Write $E_z=u+iv$ with $u,v$ real-linear. For $z\in\mathbb R$, $\fhat(z)=\int f(x)\cos(zx)\,dx$
is real; for $z=it$, $\fhat(z)=\int f(x)\cosh(tx)\,dx$ is real; in both cases $v\equiv0$,
$u\not\equiv0$, and $\Re[E_z^2]=u^2$ has signature $(1,0)$. Otherwise $(u,v):\mathcal V\to
\mathbb R^2$ is onto by Lemma~\ref{lem:indep} with $n=1$, and $u^2-v^2$ pulls back
$\mathrm{diag}(1,-1)$.
\end{proof}

\subsection{From blocks to count}
The negative index of a \emph{sum} $\sum_k(u_k^2-v_k^2)$ is not automatically the number of
blocks: the functionals could be dependent. On $\mathcal V$ they are dependent in a specific
way: $E_{-z}=E_z$ and $E_{\bar z}=\overline{E_z}$. We call $O(z)=\{\pm z,\pm\bar z\}$ the
\emph{orbit} of $z$.

\begin{lemma}[Independence of off-axis evaluations]\label{lem:indep}
Let $z_1,\dots,z_n\in\mathbb C\setminus(\mathbb R\cup i\mathbb R)$ have pairwise disjoint orbits
$O(z_i)$. Then the $2n$ functionals $\Re E_{z_i},\Im E_{z_i}$ are $\mathbb R$-linearly independent
on $\mathcal V$, the map $f\mapsto(\Re E_{z_i}(f),\Im E_{z_i}(f))_i$ is onto $\mathbb R^{2n}$, and
$\sum_{i=1}^n\big((\Re E_{z_i})^2-(\Im E_{z_i})^2\big)$ has signature $(n,n)$ on $\mathcal V$.
More generally, if in addition $y_1,\dots,y_r\in\mathbb R\cup i\mathbb R$ have pairwise disjoint
orbits $\{\pm y_j\}$, disjoint from all $O(z_i)$, then the $2n+r$ functionals $\Re E_{z_i},\Im
E_{z_i},E_{y_j}$ are linearly independent on $\mathcal V$.
\end{lemma}
\begin{proof}
Suppose $\sum_i(a_i\Re E_{z_i}+b_i\Im E_{z_i})+\sum_j e_jE_{y_j}=0$ on $\mathcal V$ with real
coefficients, and put $c_i=a_i-ib_i$, so that $a_i\Re E_{z_i}+b_i\Im E_{z_i}=\Re(c_iE_{z_i})$.
For $f\in\mathcal V$,
\[
\Re(c_iE_{z_i}(f))=\tfrac12\int f(u)\big(c_ie^{iz_iu}+\bar c_ie^{-i\bar z_iu}\big)\,du
=\int f(u)\,D_i(u)\,du,
\]
where, using that $f$ is even, $D_i(u)=\frac14\big[c_i(e^{iz_iu}+e^{-iz_iu})+\bar
c_i(e^{i\bar z_iu}+e^{-i\bar z_iu})\big]$. Similarly $E_{y_j}(f)=\int f(u)\frac12(e^{iy_ju}+
e^{-iy_ju})\,du$. Hence $\int fD=0$ for all $f\in\mathcal V$, where
$D=\sum_iD_i+\frac12\sum_je_j(e^{iy_ju}+e^{-iy_ju})$ is a smooth even real function. Testing
against all even real $f\in C_c^\infty$ gives $D=0$. The exponents $\{\pm iz_i,\pm i\bar z_i\}$
($4n$ distinct numbers, since $z_i\notin\mathbb R\cup i\mathbb R$) and $\{\pm iy_j\}$ are pairwise
distinct by the disjointness of the orbits (for $y_j=0$ the single exponent $0$ occurs, with
coefficient $e_j$). Exponentials $e^{\lambda u}$ with distinct $\lambda$ are linearly independent
functions, so every coefficient vanishes: $c_i=0$ and $e_j=0$. Independence of finitely many
functionals is equivalent to surjectivity of the joint map onto $\mathbb R^{2n+r}$; in the
coordinates $(\Re E_{z_i},\Im E_{z_i})$ the form is $\mathrm{diag}(1,-1,\dots,1,-1)$ pulled back
by a surjection, hence of signature $(n,n)$.
\end{proof}

For the zeros, the evaluation point of the quartet $\{\half\pm b\pm i\gamma\}$ is
$z=\gamma-ib$, whose orbit $\{\pm\gamma\pm ib\}$ is exactly the image of the quartet under
$\rho\mapsto(\rho-\half)/i$. Distinct quartets therefore have disjoint orbits, and $\gamma\neq0$,
$b\neq0$ give $z\notin\mathbb R\cup i\mathbb R$.

\begin{corollary}[The count]\label{cor:count}
Suppose the off-line zeros of $\xi$ form exactly $m$ quartets, $1\le m<\infty$, with evaluation
points $z_k=\gamma_k-ib_k$ and multiplicities $\mu_k$. Then:
\begin{enumerate}
\item[(a)] $Q_{\mathrm{off}}=\sum_{k=1}^m4\mu_k\big((\Re E_{z_k})^2-(\Im E_{z_k})^2\big)$ has
signature $(m,m)$ on $\mathcal V$;
\item[(b)] $\kV\le m$: indeed $Q\ge0$ on the subspace $\bigcap_k\ker(\Im E_{z_k})$, which has
codimension $m$ in $\mathcal V$.
\end{enumerate}
If there are no off-line zeros, $\kV=0$. Whether $\kV=m$ when $1\le m<\infty$ is not settled
here: a proof would require, for each prescribed vector $(\Im E_{z_k}(f))_k$, test functions
making $Q_{\mathrm{on}}(f)+\sum_k4\mu_k(\Re E_{z_k}(f))^2$ arbitrarily small, an approximation
problem of the type treated by Bombieri~\cite{Bombieri} for Weil's criterion.
\end{corollary}
\begin{proof}
(a) is Theorem~\ref{thm:loc}(2) together with Lemma~\ref{lem:indep}, the positive weights
$4\mu_k$ not affecting the signature. (b) On the stated subspace $Q=Q_{\mathrm{on}}+\sum_k
4\mu_k(\Re E_{z_k})^2\ge0$, so every subspace on which $Q$ is negative definite meets it only in
$0$ and has dimension at most $m$. If there are no off-line zeros, $Q=Q_{\mathrm{on}}\ge0$.
\end{proof}

\begin{remark}[The index $\kV$ and zeros of structure functions]\label{rem:deferred}
The number $\kV$ is the negative index of a real symmetric form on $\mathcal V$; it is not a priori
the number of zeros of an entire function in a half-plane. For a de Branges structure function $E$
write $E^\#(z)=\overline{E(\bar z)}$ and $k(E)$ for the number of zeros of $E$ in the open upper
half-plane. The function $E_\xi(z)=\xi(\half+iz)$ is \emph{not} a structure function: by the
reality of $\xi$ and the functional equation,
$E_\xi^\#(z)=\xi(\half-iz)=\xi(\half+iz)=E_\xi(z)$, so
$|E_\xi(z)|=|E_\xi(\bar z)|$ everywhere and the de Branges kernel
$[E(z)\overline{E(w)}-E^\#(z)\overline{E^\#(w)}]/(2\pi i(\bar w-z))$ vanishes identically; thus
$H(E_\xi)=\{0\}$. A genuine candidate is Lagarias' function $\xi(\half-iz)+\xi'(\half-iz)$, which
is Hermite--Biehler if and only if $\RH$ holds~\cite[Theorem~1]{Lagarias}. For structure functions with
finitely many zeros in the upper half-plane the kernel has finitely many negative squares and the
space is a Pontryagin space~\cite{KW}; for instance $E(z)=z-i$ gives the kernel $-1/\pi$, one
negative square. We do not identify $\kV$ with $k(E)$ for any specific $E$.
\end{remark}

\subsection{Signature from the evaluation points}
\begin{proposition}[Definiteness is governed by real vs.\ complex evaluation points]\label{prop:spectrum}
Let $z_1,\dots,z_n\in\mathbb C$ have pairwise disjoint orbits, $r$ of them in $\mathbb R\cup
i\mathbb R$ and $n_c=n-r$ outside it. Then the form $\sum_i\Re[E_{z_i}^2]$ on $\mathcal V$ has
signature $(r+n_c,n_c)$. In particular it is positive semidefinite if and only if all $z_i$ lie
in $\mathbb R\cup i\mathbb R$.
\end{proposition}
\begin{proof}
By Lemma~\ref{lem:indep} the functionals $E_{y}$ ($y$ in $\mathbb R\cup i\mathbb R$) and $\Re
E_{z},\Im E_{z}$ (otherwise) are jointly independent, and in these coordinates the form is the
diagonal form with $r+n_c$ entries $+1$ and $n_c$ entries $-1$.
\end{proof}

Thus, for finitely many points, indefiniteness of the bilinear evaluation form is exactly the
presence of evaluation points off $\mathbb R\cup i\mathbb R$, each giving one $(1,1)$ block
(Theorem~\ref{thm:F}). For the zeros of $\xi$ the points in $i\mathbb R$ do not occur.

\section{The index regimes and the finiteness question}\label{sec:regimes}

\begin{definition}
A \emph{Pontryagin space} $\Pi_\kappa$ is a Kre\u\i n space whose negative index $\kappa$ is
finite; a \emph{Kre\u\i n space} allows $\kappa=\infty$ (see~\cite{Bognar,Azizov}). By de
Branges~\cite{deBranges}, $H(E)$ is a Hilbert space when $E$ is Hermite--Biehler; the
corresponding Pontryagin spaces of entire functions, with finite negative index, are studied
in~\cite{KW}.
\end{definition}

\begin{proposition}[Three regimes]\label{prop:regimes}
Let $m\in\{0,1,2,\dots,\infty\}$ be the number of off-line quartets of zeros of $\xi$.
\[
\begin{array}{lll}
\text{(H)} & m=0\ (\RH): & Q\ge0\ \text{on}\ \mathcal V,\ \kV=0;\\
\text{(P)} & 0<m<\infty: & \kV\le m,\ \text{and}\ Q_{\mathrm{off}}\ \text{has signature}\ (m,m);\\
\text{(K)} & m=\infty: & \text{the arguments above give no finite bound on}\ \kV.
\end{array}
\]
\end{proposition}
\begin{proof}
(H) is Theorem~\ref{thm:loc}, (P) is Corollary~\ref{cor:count}. In case (K) the subspace used in
Corollary~\ref{cor:count}(b) has infinite codimension.
\end{proof}

\begin{proposition}[Where finiteness sits]\label{prop:finiteness}
The statement $m<\infty$ (finitely many off-line zeros) is implied by $\RH$ ($m=0$) and is a
priori weaker; it is not implied by the known unconditional results. The positive-proportion
theorems give a positive proportion of zeros on the line (at least $1/3$ by
Levinson~\cite{Levinson}, at least $2/5$ by Conrey~\cite{Conrey}, at least $5/12$ by Pratt,
Robles, Zaharescu and Zeindler~\cite{PRZZ}) but bound nothing off it. The zero-density estimates,
for instance Selberg's $N(\sigma,T)\ll T^{1-(\sigma-\frac12)/4}\log T$~\cite{Selberg}, give
$o(T)$ zeros with $\beta>\sigma$ up to height $T$ for each fixed $\sigma>\half$, but permit
infinitely many in total. By Corollary~\ref{cor:count}, each quartet contributes at most one
negative direction of $Q$ on $\mathcal V$; this does not bound $m$.
\end{proposition}

\begin{remark}
Thus the finite-versus-infinite distinction for $\kV$ is tied to the zero-density problem. The
analysis of the regime (K) would require tools beyond those used here. For the full Weil form
$A_\infty-P$ restricted to functions supported in a fixed interval, finiteness of the negative
index is proved in~\cite{P11}. By Remark~\ref{rem:choices}, $Q(f,f)$ is the value of that form at
$f$ when $f\in\mathcal V$ is supported in $[-d,d]$; hence the negative index of $Q$ on
$\mathcal V_d=\{f\in\mathcal V:\operatorname{supp}f\subset[-d,d]\}$ is finite for every $d$, whatever
$m$ is, and $\kV$ is the supremum of these indices over $d$. This bounds neither $\kV$ nor $m$.
\end{remark}

\section{The de Branges axioms for a realization}\label{sec:rigidity}

We ask what the axiomatic characterization of de Branges spaces, extended to Pontryagin spaces by
Kaltenb\"ack and Woracek~\cite{KW}, says about a space carrying $Q$.

\begin{definition}[de Branges axioms; realizations]\label{def:axioms}
Let $\Pi$ be a Pontryagin space of entire functions with indefinite inner product
$[\cdot,\cdot]$, linear in the first and conjugate-linear in the second argument; write
$\mathcal E_w:f\mapsto f(w)$ and $f^\#(z)=\overline{f(\bar z)}$. Then $\Pi$ is a
\emph{de Branges--Pontryagin space} if:
\begin{itemize}
\item[\textbf{(dB1)}] $\mathcal E_w$ is continuous for every non-real $w$;
\item[\textbf{(dB2)}] \emph{(difference-quotient axiom)} if $f\in\Pi$ and $f(w)=0$ for a
non-real $w$, then $z\mapsto f(z)(z-\bar w)/(z-w)$ belongs to $\Pi$ and has the same
$[\cdot,\cdot]$-norm as $f$;
\item[\textbf{(dB3)}] \emph{(conjugation isometry)} $f\mapsto f^\#$ maps $\Pi$ onto $\Pi$ with
$[f^\#,g^\#]=[g,f]$.
\end{itemize}
A \emph{realization} of $Q$ is a Pontryagin space $\Pi$ of entire functions together with a
real-linear map $J:\mathcal V\to\Pi$ such that $[Jf,Jg]=Q(f,g)$ for all $f,g\in\mathcal V$ and the
complex linear span of $J(\mathcal V)$ is dense in $\Pi$.
\end{definition}

If $f/(z-w)$ and $g/(z-v)$ belong to $\Pi$ and multiplication by $z$ is symmetric on these
elements, then for $f(w)=g(v)=0$,
\[
\Big[\tfrac{f(z)}{z-w},g\Big]-\Big[f,\tfrac{g(z)}{z-v}\Big]
=(w-\bar v)\Big[\tfrac{f(z)}{z-w},\tfrac{g(z)}{z-v}\Big],
\]
as one sees by writing $g=(z-v)\frac{g}{z-v}$ and $f=(z-w)\frac{f}{z-w}$; this identity is the
form in which (dB2) is often used.

\begin{theorem}[Kaltenb\"ack--Woracek~\cite{KW}]\label{thm:dB}
A Pontryagin space of entire functions with negative index $k<\infty$, satisfying
\textnormal{(dB1)--(dB3)} and containing a nonzero element, equals isometrically a space $H(E)$
for an entire function $E$ with exactly $k$ zeros in the open upper half-plane (counted with
multiplicity, under the normalizations of~\cite{KW}). The function $E=A-iB$ is determined by the
space up to replacing $(A,B)$ by $(A,B)M$ with $M\in SL_2(\mathbb R)$. For $k=0$ this is de
Branges' theorem~\cite{deBranges}.
\end{theorem}

\begin{proposition}[Cost of each axiom]\label{prop:auto}
Let $(\Pi,J)$ be a realization of $Q$.
\begin{enumerate}
\item If $(Jf)^\#=Jf$ for every $f\in\mathcal V$ \textnormal{(}as when $Jf$ is built from
$\varphi_f$, since $\varphi_f(\bar s)=\overline{\varphi_f(s)}$\textnormal{)} and $f\mapsto
f^\#$ is continuous on $\Pi$, then \textnormal{(dB3)} holds.
\item \textnormal{(dB1)} is not a consequence of the definition: Theorem~\ref{thm:loc} concerns only
the functionals $E_{\gamma-ib}$ at the zeros, whereas \textnormal{(dB1)} demands continuity of
point evaluation at every non-real $w$ of the functions in $\Pi$.
\item \textnormal{(dB2)} is not a consequence of the symmetries of $Q$: it is a closure property of
$\Pi$ under a division operation.
\end{enumerate}
\end{proposition}
\begin{proof}
(1) Let $a,b,c,d\in J(\mathcal V)$. Since $[a,c]=Q(\cdot,\cdot)$ is real and symmetric on
$J(\mathcal V)$, and $a^\#=a$ etc.,
\[
[(a+ib)^\#,(c+id)^\#]=[a-ib,c-id]=[a,c]+i[a,d]-i[b,c]+[b,d],
\]
while $[c+id,a+ib]=[c,a]-i[c,b]+i[d,a]+[d,b]$; these agree. So $f\mapsto f^\#$ is a
conjugate-linear map of the dense span of $J(\mathcal V)$ satisfying $[f^\#,g^\#]=[g,f]$; by
continuity this extends to $\Pi$, and $f^{\#\#}=f$ gives surjectivity. (2)--(3) These are
statements about all elements of $\Pi$, not determined by the values of $Q$ on $\mathcal V$.
\end{proof}

\begin{theorem}[Structure of a de Branges realization]\label{thm:rigid}
Let $(\Pi,J)$ be a realization of $Q$ with finite negative index $k$, satisfying
\textnormal{(dB1)--(dB3)}, with $\Pi\neq\{0\}$. Then:
\begin{enumerate}
\item $\Pi=H(E)$ isometrically for some entire $E$ with exactly $k$ zeros in the open upper
half-plane;
\item $k\ge\kV$.
\end{enumerate}
\end{theorem}
\begin{proof}
(1) is Theorem~\ref{thm:dB}. (2) Let $N\subset\mathcal V$ be a subspace of dimension $d$ on which
$Q$ is negative definite, with basis $n_1,\dots,n_d$. For complex $c$,
$[\sum c_iJn_i,\sum c_jJn_j]=\sum_{i,j}c_i\bar c_jQ(n_i,n_j)$, and a real symmetric negative
definite matrix is negative definite as a Hermitian matrix. Hence the complex span of
$Jn_1,\dots,Jn_d$ is a $d$-dimensional negative definite subspace of $\Pi$ (the $Jn_i$ are
independent since the Gram matrix is invertible), and $d\le k$.
\end{proof}

\begin{remark}[Location of the zeros of $E$]\label{rem:identify}
Theorem~\ref{thm:rigid} does not say where the $k$ zeros of $E$ in the upper half-plane lie, and
in particular does not relate them to the off-line zeros of $\xi$: equal negative indices do not
force equal zero sets. The function $E_\xi$ cannot serve as $E$, since $H(E_\xi)=\{0\}$
(Remark~\ref{rem:deferred}). Whether a realization exists whose structure function has its
upper-half-plane zeros at the points $z_\rho$ of the off-line zeros is an open question.
\end{remark}

\begin{corollary}[The open items]\label{cor:freedom}
For a realization $(\Pi,J)$ of $Q$ with finite negative index and $(Jf)^\#=Jf$, the conjugation
axiom holds under the continuity condition of Proposition~\ref{prop:auto}(1), and the remaining
conditions for Theorem~\ref{thm:rigid} are \textnormal{(i)} continuity of point evaluations
\textnormal{(dB1)} and \textnormal{(ii)} the difference-quotient axiom \textnormal{(dB2)}. Neither
is supplied by the symmetries of $Q$, and even when both hold the relation between the zeros of
$E$ and those of $\xi$ remains open. If $\kV=\infty$ there is no realization in a Pontryagin
space, by the argument of Theorem~\ref{thm:rigid}(2).
\end{corollary}

\section{Symmetry versus positivity}\label{sec:dichotomy}

Write $\sigma(s)=\bar s$, $\iota(s)=1-s$, and $j=\sigma\iota:\,s\mapsto1-\bar s$, an
anti-holomorphic involution with $\Fix(j)=\{\Re s=\half\}$.

\begin{definition}[Symmetry data]\label{def:real}
The \emph{symmetry data} of the zeros is the pair consisting of the zero multiset
$\{\rho\}$ and the action of the group generated by $\sigma$ and $\iota$ on it. Its
\emph{positivity} is the condition $Q\ge0$ on $\mathcal V$, implied by $\RH$
(Theorem~\ref{thm:loc}).
\end{definition}

\begin{proposition}[Symmetry is unconditional]\label{prop:sym}
$\sigma$ and $\iota$ generate a Klein four-group $V_4\cong(\mathbb Z/2)^2$ acting on $\mathbb C$;
$j=\sigma\iota$ is an anti-holomorphic involution with $\Fix(j)$ the critical line; and the zero
set of $\xi$, with multiplicities, is $V_4$-invariant.
\end{proposition}
\begin{proof}
$\sigma$ and $\iota$ are commuting involutions ($\sigma\iota(s)=1-\bar s=\iota\sigma(s)$),
$j^2=\mathrm{id}$, and $1-\bar s=s\iff\Re s=\half$. The functional equation and the reality of
$\xi$ give $V_4$-invariance of the zeros with multiplicities.
\end{proof}

Whether every admissible configuration of zeros admits a realization in the sense of
Definition~\ref{def:axioms} satisfying \textnormal{(dB1)--(dB3)} is a question we leave open.

\begin{proposition}[Setwise and pointwise invariance]\label{prop:kappa}
The involution $j$ acts on the zeros setwise and possibly pointwise:
\[
\text{(setwise)}\ \ j(\{\rho\})=\{\rho\};\qquad
\text{(pointwise)}\ \ j(\rho)=\rho\ \forall\rho .
\]
The setwise statement holds unconditionally; the pointwise statement is $\{\rho\}\subset\Fix(j)$,
i.e.\ $\RH$. The form $Q$ is unchanged if the summation is relabelled by any element of $V_4$.
\end{proposition}
\begin{proof}
An off-line $\rho$ has the $j$-invariant orbit $\{\half\pm b\pm i\gamma\}$, so $j$ permutes the
zeros regardless of $\RH$; $j$ fixes $\rho$ iff $\Re\rho=\half$. The last statement is
Lemma~\ref{lem:conv}: the summand $\varphi_f(\rho)\varphi_g(\rho)$ is invariant under $\iota$ and
conjugated under $\sigma$, and $Q$ is real.
\end{proof}

\begin{remark}[Heuristic: negative directions are not local]\label{prop:local}
The following is a heuristic description, not a theorem. Suppose $\rho_0=\half+b+i\gamma_0$ is an
off-line zero. The negative direction of the block $4(u^2-v^2)$ asks for $f$ with
$\Re\fhat(\gamma_0-ib)$ small and $\Im\fhat(\gamma_0-ib)$ large. If $f$ has support of length
$L$, then $\fhat$ is spread over frequencies of width about $1/L$ around $\gamma_0$, which can be
smaller than the mean spacing $2\pi/\log(\gamma_0/2\pi)$ of the zeros; one then expects the
on-line contribution $\sum\fhat(\gamma)^2$ from neighbouring zeros to be controlled only through
the global identity \eqref{eq:weil-id}, $Q(f,f)=A_\infty(f*f)-P(f*f)$, at the scale of
Montgomery's pair correlation~\cite{Montgomery}. We do not make this precise.
\end{remark}

\begin{theorem}[Symmetry versus positivity]\label{thm:dichotomy}
\emph{(i)} The symmetry data of Definition~\ref{def:real} is present unconditionally: the zero set
is $V_4$-invariant (Proposition~\ref{prop:sym}). \emph{(ii)} $\RH$ is the pointwise condition
$\{\rho\}\subset\Fix(j)$, which is not implied by the setwise invariance
(Proposition~\ref{prop:kappa}); equivalently, $m=0$ in the notation of
Proposition~\ref{prop:regimes}, and it implies $Q\ge0$ on $\mathcal V$.
\end{theorem}
\begin{proof}
(i) is Proposition~\ref{prop:sym}. For (ii), $V_4$-invariant configurations with off-line quartets
are consistent with the group action (the action on a quartet is free), so setwise invariance
does not force $m=0$; the identification of $\RH$ with $\{\rho\}\subset\Fix(j)$ is
Proposition~\ref{prop:kappa}, and $Q\ge0$ under $\RH$ is Theorem~\ref{thm:loc}.
\end{proof}

\begin{remark}[Heuristic: spectral programmes]\label{cor:converge}
The Hilbert--P\'olya / de Branges canonical-system approach, Connes' trace
formula~\cite{Connes99} and Burnol's de Branges spaces~\cite{Burnol} each incorporate the
$V_4$-symmetry of the zeros in their construction. Theorem~\ref{thm:dichotomy} shows only that
this symmetry alone does not imply $\RH$; we make no claim about what further structure these
approaches provide.
\end{remark}

\section{Conclusion}

The proved results are as follows. On even real test functions the localized Weil form splits into
nonnegative on-line squares and one Lorentzian block per off-line quartet
(Theorem~\ref{thm:loc}); evaluation functionals at points with disjoint orbits are independent
(Lemma~\ref{lem:indep}), so with $m<\infty$ off-line quartets the off-line part has signature
$(m,m)$ and the negative index satisfies $\kV\le m$ (Corollary~\ref{cor:count}). Forms realized by
positive kernels, one-sided Cauchy Gram forms and continuous non-degenerate families on
$\mathbb C^N$ cannot exhibit negative directions that appear exactly off the line
(Theorems~\ref{thm:A}--\ref{thm:C}, Corollary~\ref{cor:class}); the modular resonance pairings are
an instance (Theorem~\ref{thm:modular}), and their weights are computed numerically and
asymptotically (Proposition~\ref{prop:weights}). A realization of $Q$ in a Pontryagin space of
entire functions satisfying the three de Branges axioms is a space $H(E)$ with $k(E)\ge\kV$
(Theorem~\ref{thm:rigid}), and $E_\xi(z)=\xi(\half+iz)$ is not a structure function
(Remark~\ref{rem:deferred}).

Open are: the equality $\kV=m$ for finitely many off-line quartets; the finiteness of $m$; the
existence of a realization satisfying (dB1)--(dB2); and the location of the zeros of its
structure function. Nothing here proves the Riemann Hypothesis.

\begin{thebibliography}{99}
\bibitem{P7} D.~A.~Trejo Pizzo, \emph{The localized Weil form on a Hermite--Gauss basis: the
effect of an off-line zero pair and the truncation of the prime sum}, Paper~7 of this series.
\bibitem{P8} D.~A.~Trejo Pizzo, \emph{A Kre\u{\i}n-space formulation of the zero side of the Weil
quadratic form on Paley--Wiener spaces}, Paper~8 of this series.
\bibitem{P11} D.~A.~Trejo Pizzo, \emph{The off-line zeros as a Pontryagin index: a conditional
reduction and numerical observations on the prime tail}, Paper~11 of this series.
\bibitem{Aronszajn} N.~Aronszajn, \emph{Theory of reproducing kernels}, Trans.\ Amer.\ Math.\
Soc.\ \textbf{68} (1950), 337--404.
\bibitem{Azizov} T.~Ya.~Azizov and I.~S.~Iokhvidov, \emph{Linear Operators in Spaces with an
Indefinite Metric}, Wiley, 1989.
\bibitem{Bognar} J.~Bogn\'ar, \emph{Indefinite Inner Product Spaces}, Springer, 1974.
\bibitem{Bombieri} E.~Bombieri, \emph{Remarks on Weil's quadratic functional in the theory of
prime numbers, I}, Rend.\ Mat.\ Acc.\ Lincei (9) \textbf{11} (2000), 183--233.
\bibitem{Burnol} J.-F.~Burnol, \emph{Sur les ``espaces de Sonine'' associ\'es par de Branges \`a
la transformation de Fourier}, C.~R.\ Math.\ Acad.\ Sci.\ Paris \textbf{335} (2002), 689--692.
\bibitem{Connes99} A.~Connes, \emph{Trace formula in noncommutative geometry and the zeros of the
Riemann zeta function}, Selecta Math.\ (N.S.) \textbf{5} (1999), 29--106.
\bibitem{Conrey} J.~B.~Conrey, \emph{More than two fifths of the zeros of the Riemann zeta
function are on the critical line}, J.\ Reine Angew.\ Math.\ \textbf{399} (1989), 1--26.
\bibitem{deBranges} L.~de Branges, \emph{Hilbert Spaces of Entire Functions}, Prentice--Hall,
Englewood Cliffs NJ, 1968.
\bibitem{Iwaniec} H.~Iwaniec, \emph{Spectral Methods of Automorphic Forms}, 2nd ed., Grad.\ Stud.\
Math.\ \textbf{53}, Amer.\ Math.\ Soc., 2002.
\bibitem{KW} M.~Kaltenb\"ack and H.~Woracek, \emph{Pontryagin spaces of entire functions I},
Integral Equations Operator Theory \textbf{33} (1999), 34--97.
\bibitem{Lagarias} J.~C.~Lagarias, \emph{Hilbert spaces of entire functions and Dirichlet
$L$-functions}, in: Frontiers in Number Theory, Physics, and Geometry~I, Springer, Berlin, 2006,
365--377.
\bibitem{Levinson} N.~Levinson, \emph{More than one third of zeros of Riemann's zeta-function are
on $\sigma=1/2$}, Adv.\ Math.\ \textbf{13} (1974), 383--436.
\bibitem{Montgomery} H.~L.~Montgomery, \emph{The pair correlation of zeros of the zeta
function}, Proc.\ Sympos.\ Pure Math.\ \textbf{24} (1973), 181--193.
\bibitem{PRZZ} K.~Pratt, N.~Robles, A.~Zaharescu and D.~Zeindler, \emph{More than five-twelfths of
the zeros of $\zeta$ are on the critical line}, Res.\ Math.\ Sci.\ \textbf{7} (2020), no.~1, Paper No.~2.
\bibitem{Selberg} A.~Selberg, \emph{Contributions to the theory of the Riemann zeta-function},
Arch.\ Math.\ Naturvid.\ \textbf{48} (1946), 89--155.
\bibitem{Weil} A.~Weil, \emph{Sur les ``formules explicites'' de la th\'eorie des nombres
premiers}, Comm.\ S\'em.\ Math.\ Univ.\ Lund (1952), 252--265.
\bibitem{Zagier} D.~Zagier, \emph{The Rankin--Selberg method for automorphic functions which
are not of rapid decay}, J.\ Fac.\ Sci.\ Univ.\ Tokyo \textbf{28} (1981), 415--437.
\end{thebibliography}

\end{document}
